\section{Performance-Guaranteed Dimension Reduction}
\label{sec:dimension-reduction}

The theoretical bound is used before the downstream scheduling MILP is constructed. The contribution is a guaranteed preprocessing framework rather than a new MILP solution algorithm.

\subsection{Scheduling dimension}

If the vehicle-level formulation introduces one precedence variable for each pair of vehicles from conflicting approaches, its pairwise ordering dimension is
\begin{equation}
C_0=\sum_{l<l'}n_ln_{l'}.
\end{equation}
After platooning, the corresponding dimension becomes
\begin{equation}
C(\Pi)=\sum_{l<l'}K_lK_{l'}.
\label{eq:dimension-measure}
\end{equation}
This quantity is a deterministic measure of model size. Actual Gurobi runtime is evaluated empirically and is not assumed to be a strictly monotone function of \(C(\Pi)\).

\subsection{Loss-budget formulation}

Given an acceptable average-delay loss \(\varepsilon\), the central scheduler selects a partition by solving
\begin{align}
\min_{\Pi}\quad & C(\Pi)\\
\text{s.t.}\quad & B(\Pi)\le\varepsilon,\\
& \Pi\text{ is a contiguous partition},\\
& |P_{l,j}|\le P_{\max},\quad \forall l,j.
\label{eq:loss-budget-formulation}
\end{align}
The downstream scheduling MILP is then built using the selected platoons and solved with Gurobi.

\subsection{Computation-size formulation}

Alternatively, if the scheduler can process at most \(\overline C\) pairwise ordering decisions, it solves
\begin{align}
\min_{\Pi}\quad & B(\Pi)\\
\text{s.t.}\quad & C(\Pi)\le\overline C,\\
& \Pi\text{ is a contiguous partition},\\
& |P_{l,j}|\le P_{\max},\quad \forall l,j.
\label{eq:size-budget-formulation}
\end{align}
The two formulations generate a set of nondominated dimension--loss pairs
\begin{equation}
\mathcal{F}
=
\left\{
\big(C(\Pi),B(\Pi)\big):
\Pi\text{ is feasible}
\right\}.
\end{equation}

No specialized solution algorithm is claimed. The partition model and downstream scheduling model are both solved by an off-the-shelf optimizer. The methodological contribution lies in the valid loss guarantee and its use in selecting the aggregation level.
